\documentclass[12pt]{article}
\usepackage{amsmath, amssymb}
\usepackage{geometry}
\geometry{margin=1in}

\title{Lecture Scribe: Transformation of Random Variables and Functions of Two Random Variables}
\author{}
\date{}

\begin{document}

\maketitle

\section{Fundamental Framework}

\subsection{Given}

A continuous random variable (CRV):
\[
X
\]

whose probability functions are known:

\textbf{Probability Density Function (PDF):}
\[
f_X(x)
\]

\textbf{Cumulative Distribution Function (CDF):}
\[
F_X(x)
\]

This is known \textbf{a priori}, meaning all probability information about $X$ is available.

\subsection{Definition: Transformation of Random Variable}

A new random variable is defined as:

\[
Y = g(X)
\]

This produces a new random variable derived from $X$.

\subsection{Objective}

Determine:

\[
F_Y(y)
\quad \text{and} \quad
f_Y(y)
\]

Logical dependency: Since $Y$ is defined in terms of $X$, its distribution must be derived from the distribution of $X$.

\section{Invertibility Requirement}

Assume:

\[
g(x)
\]

is monotonic (either increasing or decreasing).

This ensures existence of inverse:

\[
x = g^{-1}(y)
\]

This inverse relationship allows conversion of probabilities involving $Y$ into probabilities involving $X$.

\section{CDF Method for Transformation}

\subsection{Step 1: Definition of CDF}

By definition:

\[
F_Y(y) = Pr(Y \le y)
\]

Substitute $Y = g(X)$:

\[
F_Y(y) = Pr(g(X) \le y)
\]

\section{Case 1: Monotonically Increasing Function}

\subsection{Step 2: Inequality Transformation}

Since $g(x)$ is increasing:

\[
g(X) \le y \iff X \le g^{-1}(y)
\]

\subsection{Step 3: Express Using Known CDF}

\[
F_Y(y)
=
Pr(X \le g^{-1}(y))
\]

Using definition of $F_X(x)$:

\[
\boxed{
F_Y(y) = F_X(g^{-1}(y))
}
\]

\subsection{Step 4: Differentiate to Obtain PDF}

\[
f_Y(y)
=
\frac{d}{dy}F_Y(y)
\]

\[
=
\frac{d}{dy}F_X(g^{-1}(y))
\]

Using chain rule:

\[
=
f_X(g^{-1}(y))
\frac{d}{dy}g^{-1}(y)
\]

Since $x = g^{-1}(y)$:

\[
\boxed{
f_Y(y)
=
f_X(x)
\left|
\frac{dx}{dy}
\right|
}
\]

\section{Case 2: Monotonically Decreasing Function}

\subsection{Step 1: Define CDF}

\[
F_Y(y)
=
Pr(g(X) \le y)
\]

Since decreasing:

\[
g(X) \le y \iff X \ge g^{-1}(y)
\]

\subsection{Step 2: Use Complement Rule}

\[
F_Y(y)
=
Pr(X \ge g^{-1}(y))
\]

\[
=
1 - F_X(g^{-1}(y))
\]

\subsection{Step 3: Differentiate}

\[
f_Y(y)
=
-
f_X(g^{-1}(y))
\frac{d}{dy}g^{-1}(y)
\]

Final form:

\[
\boxed{
f_Y(y)
=
\frac{f_X(x)}
{\left|\frac{dy}{dx}\right|}
}
\]

\section{Example: Uniform Transformation}

\subsection{Given}

\[
X \sim Uniform(-1,1)
\]

\[
f_X(x) =
\begin{cases}
\frac{1}{2}, & -1<x<1 \\
0, & otherwise
\end{cases}
\]

Transformation:

\[
Y =
\sin\left(\frac{\pi X}{2}\right)
\]

\subsection{Step 1: Inverse}

\[
x =
\frac{2}{\pi}
\sin^{-1}(y)
\]

\subsection{Step 2: Derivative}

\[
\frac{dx}{dy}
=
\frac{2}{\pi}
\frac{1}{\sqrt{1-y^2}}
\]

\subsection{Step 3: Apply Formula}

\[
f_Y(y)
=
f_X(x)
\left|\frac{dx}{dy}\right|
\]

\[
=
\frac{1}{2}
\cdot
\frac{2}{\pi}
\cdot
\frac{1}{\sqrt{1-y^2}}
\]

Final result:

\[
\boxed{
f_Y(y)
=
\frac{1}{\pi \sqrt{1-y^2}},
\quad -1<y<1
}
\]

\section{Function of Two Random Variables}

Define:

\[
Z = g(X,Y)
\]

Example:

\[
Z = X + Y
\]

Goal:

\[
f_Z(z)
\]

Requires joint PDF:

\[
f_{XY}(x,y)
\]

\section{CDF of Z}

\[
F_Z(z)
=
Pr(Z \le z)
\]

\[
=
Pr(X+Y \le z)
\]

Region:

\[
x+y \le z
\]

Integral form:

\[
F_Z(z)
=
\int_{-\infty}^{\infty}
\int_{-\infty}^{z-y}
f_{XY}(x,y)
dx dy
\]

\section{Leibnitz Rule}

Given:

\[
G(x)
=
\int_{a(x)}^{b(x)}
h(x,y)
dy
\]

Derivative:

\[
\frac{dG(x)}{dx}
=
\frac{db(x)}{dx}h(x,b(x))
-
\frac{da(x)}{dx}h(x,a(x))
+
\int_{a(x)}^{b(x)}
\frac{\partial h(x,y)}{\partial x}
dy
\]

\section{PDF of Z}

Differentiate CDF:

\[
\boxed{
f_Z(z)
=
\int_{-\infty}^{\infty}
f_{XY}(z-y,y)
dy
}
\]

Equivalent form:

\[
\boxed{
f_Z(z)
=
\int_{-\infty}^{\infty}
f_{XY}(x,z-x)
dx
}
\]

\section{Independent Case}

If independent:

\[
f_{XY}(x,y)
=
f_X(x)f_Y(y)
\]

\[
\boxed{
f_Z(z)
=
\int_{-\infty}^{\infty}
f_X(x)f_Y(z-x)dx
}
\]

This is convolution.

\section{Normal Example}

\[
X \sim N(0,1)
\]

\[
Y \sim N(0,1)
\]

\[
\boxed{
Z \sim N(0,2)
}
\]

\section{Exponential Example}

\[
f_X(x) = \lambda e^{-\lambda x}
\]

\[
f_Y(y) = \lambda e^{-\lambda y}
\]

\[
\boxed{
f_Z(z)
=
\lambda^2 z e^{-\lambda z}
}
\]

\section{Difference of Random Variables}

\[
Z = X - Y
\]

\[
F_Z(z)
=
Pr(X-Y \le z)
\]

\[
=
Pr(X \le z+Y)
\]

Integral form:

\[
F_Z(z)
=
\int
\int_{x \le z+y}
f_{XY}(x,y)
dx dy
\]

\section{Ratio of Random Variables}

\[
Z =
\frac{X}{Y}
\]

\[
F_Z(z)
=
Pr\left(\frac{X}{Y} \le z\right)
\]

Requires cases:

\[
Y > 0, \quad Y < 0
\]

\section{Radius Transformation}

\[
R =
\sqrt{X^2 + Y^2}
\]

Goal:

\[
f_R(r)
\]

Requires transformation of joint distribution.

\section{Key Final Formulas}

\[
\boxed{
f_Y(y)
=
f_X(g^{-1}(y))
\left|
\frac{d}{dy}g^{-1}(y)
\right|
}
\]

\[
\boxed{
f_Z(z)
=
\int f_X(x)f_Y(z-x)dx
}
\]

\[
\boxed{
N(0,1)+N(0,1)=N(0,2)
}
\]

\[
\boxed{
f_Z(z)=\lambda^2 z e^{-\lambda z}
}
\]

\end{document}
